\section{A nonsymmetric periodic realization with four measured channels}
\label{sec:v44-realization}

We give a family in which the contact geometry, the two deck cycles and
intrinsic registration can all be evaluated in one construction.  The
parameters used to construct the table are not supplied to its observer.
In particular, the coordinates in the following construction verify
realizability; they are not additional observations.  The example uses
selected channels, not a decomposition of the whole collision-count event.

Put $n(\theta)=(\cos\theta,\sin\theta)$ and
$t(\theta)=(-\sin\theta,\cos\theta)$.  If $h+h''>0$, the curve
\begin{equation}\label{eq:v44-support-parametrization}
 x_h(\theta)=h(\theta)n(\theta)+h'(\theta)t(\theta)
\end{equation}
has derivative $(h+h'')t$, outward normal $n$, support function $h$,
and curvature $(h+h'')^{-1}$.  To see the support assertion directly,
for fixed $\theta_0$ the derivative of
$x_h(\theta)\cdot n(\theta_0)$ is
$(h+h'')(\theta)\sin(\theta_0-\theta)$.  Its unique maximum occurs at
$\theta=\theta_0$.  Thus \eqref{eq:v44-support-parametrization} is the
boundary of a strictly convex body, denoted by $K_h$.

Consider the two support functions
\begin{equation}\label{eq:v44-supports}
 \begin{split}
 h_1(\theta)&=1+\epsilon\left(\cos2\theta+
                  \tfrac14\cos(3\theta+\pi/4)\right),\\
 h_2(\theta)&=\tfrac34+\delta\left(\cos2\theta+
                  \tfrac14\cos(3\theta+\pi/6)\right),
 \qquad \epsilon,\delta\in[1/200,1/100].
 \end{split}
\end{equation}
Let $T$ be a real $2\times2$ matrix with $\|T-I\|\le1/100$, and let
$|w|\le1/100$.  Define
\begin{equation}\label{eq:v44-lattice-family}
 L=T\begin{pmatrix}8&0\\0&10\end{pmatrix},\qquad
 v=T(4,5)^t+w,\qquad
 D_1=K_{h_1},\quad D_2=v+K_{h_2}.
\end{equation}
The periodic obstacles are $D_a+L m$, $a\in\{1,2\}$, $m\in\mathbb Z^2$.
The marked basis of the abstract deck group is denoted by
$\mathsf e_1,\mathsf e_2$; neither column of $L$ is given as a datum.

\begin{theorem}[A uniformly admissible nonsymmetric family]
\label{thm:v44-realized-family}
Every member of \eqref{eq:v44-supports}--\eqref{eq:v44-lattice-family}
is a periodic dispersing table with connected analytic boundaries.
Every two distinct obstacle closures have distance greater than $43/10$.
The curvature ranges on the two obstacle types are disjoint:
\begin{equation}\label{eq:v44-curvature-bounds}
 \frac{20}{21}\le\kappa_1\le\frac{20}{19},\qquad
 \frac54\le\kappa_2\le\frac{10}{7}.
\end{equation}
Neither obstacle has a nonidentity Euclidean symmetry.  Modulo its period
translations, the complete table has no Euclidean symmetry either.

The four labelled channels
\begin{equation}\label{eq:v44-channel-labels}
 e_{ij}=(1,2,-i\mathsf e_1-j\mathsf e_2),\qquad i,j\in\{0,1\},
\end{equation}
are unobstructed closest-pair channels for their specified obstacle pairs.
Every such closest segment has distance at least $383/200$ from every
third obstacle.  Their contacts and gaps depend real analytically on the
parameters and have common contact and positive-offset collars.
All their incidence transitions are signature-rigid.  The cycles
\begin{equation}\label{eq:v44-cycles}
 c_1=e_{00}e_{10}^{-1},\qquad c_2=e_{00}e_{01}^{-1}
\end{equation}
form a rank-two anchoring pair based at $e_{00}$; the edge $e_{00}$ is a
signature-rigid spanning tree rooted at obstacle $1$.
\end{theorem}

\begin{proof}
For $h=r+z(\cos2\theta+\tfrac14\cos(3\theta+\phi))$,
\[
 h+h''=r-3z\cos2\theta-2z\cos(3\theta+\phi).
\]
The curvature radii lie in $[19/20,21/20]$ and $[7/10,4/5]$,
respectively.  This proves strict convexity and
\eqref{eq:v44-curvature-bounds}.  The support bound gives
$K_{h_a}\subset\overline B(0,R)$ with $R=81/80$ for both types.
At $T=I,w=0$, the set of obstacle centers is
\[
 \{(8m,10n):m,n\in\mathbb Z\}\ \cup\
 \{(4+8m,5+10n):m,n\in\mathbb Z\}.
\]
Its minimum separation is $\sqrt{41}$.  Since
$|Tz|\ge(99/100)|z|$, all distinct perturbed centers have separation
at least $(99/100)\sqrt{41}-1/100$.  Subtracting $2R$ gives
\[
 \operatorname{dist}(D_a+Lm,D_b+Ln)
 \ge\tfrac{99}{100}\sqrt{41}-\tfrac{407}{200}
 >\tfrac{43}{10}
\]
for distinct obstacles; the last strict inequality follows already from
$\sqrt{41}>32/5$.  The smallest singular value of $T$ is positive and
$\det T>0$, so the quotient is a genuine flat torus.  Its free area is
positive.  More explicitly, integration of
$\tfrac12h(h+h'')$ gives
\begin{equation}\label{eq:v44-free-area}
 A=80\det T-\frac{25\pi}{16}
                   +\frac{7\pi}{4}(\epsilon^2+\delta^2)>0.
\end{equation}

We next exclude symmetry without imposing any contact registration.
The first Fourier coefficients of both supports vanish.  The vector
\begin{equation}\label{eq:v44-center}
 s(K)=\frac1\pi\int_0^{2\pi}h_K(\theta)n(\theta)\,d\theta
\end{equation}
is covariant under translations and orthogonal transformations, so any
symmetry of either body fixes the origin.  A rotational symmetry through
$\alpha$ would preserve its nonzero modes $2$ and $3$, and hence satisfy
$e^{2i\alpha}=e^{3i\alpha}=1$.  Thus it is the identity.  Reflection in
a line at angle $\alpha$ would require
$2\alpha\in\pi\mathbb Z$ and $3\alpha+\phi\in\pi\mathbb Z$.
Neither $\phi=\pi/4$ nor $\phi=\pi/6$ permits this: substituting
$\alpha=k\pi/2$ gives, respectively, $(6k+1)/4$ and $(9k+1)/6$ as
nonintegral multiples of $\pi$.  The two perimeters, $2\pi$ and
$3\pi/2$, differ.  A symmetry of the periodic table must therefore
preserve obstacle types.  Its orthogonal part would be a symmetry of
$K_{h_1}$ and is the identity; its translation belongs to $L\mathbb Z^2$.

For clearance, first consider the center segment from $(0,0)$ to $(4,5)$.
Every third reference center has distance at least $4$ from the rectangle
$[0,4]\times[0,5]$: this follows directly by separating the integer
coordinates of the two displayed center sets.  The other three rectangles
are reflections of this one.  After applying $T$, the distance of a third
center from the corresponding center segment is at least $4(99/100)$.
Moving the second-type centers by $w$ changes a segment by at most $|w|$
in Hausdorff distance and a third center by at most $|w|$.  Thus the
lower bound becomes $4(99/100)-2/100$.  A segment joining points of the
two radius-$R$ bodies is contained in the radius-$R$ neighborhood of the
center segment, by convex interpolation of its endpoints.  Subtracting a
further $R$ for the third body proves the claimed clearance:
\begin{equation}\label{eq:v44-clearance}
 4\frac{99}{100}-\frac2{100}-2\frac{81}{80}
                         =\frac{383}{200}.
\end{equation}
This argument covers every lattice translate, not a finite numerical list.

For completeness, the contacts can be found from a scalar equation.  Put
\[
 v_{ij}=T(4-8i,5-10j)^t+w,\qquad
 d_{ij}=|v_{ij}|,\quad \vartheta_{ij}=\arg v_{ij},
\]
and maximize
\begin{equation}\label{eq:v44-gap-function}
 F_{ij}(\theta)=v_{ij}\cdot n(\theta)-h_1(\theta)-h_2(\theta+\pi).
\end{equation}
The variable $d_{ij}$ here is a center distance, not the positive time
offset $d$ used elsewhere.  The perturbation in $h_1+h_2$ from $7/4$
is at most $1/40$.  Comparison with $\theta=\vartheta_{ij}$ excludes
$|\theta-\vartheta_{ij}|\ge\pi/3$ from the set of maxima.  On the
remaining arc,
\[
 F_{ij}''\le-d_{ij}/2+1/8<0,
\]
because $\|h_1''\|+\|h_2''\|\le1/8$.  There is a unique maximum
$\theta_{ij}$ and it depends analytically on the parameters.  With
\begin{equation}\label{eq:v44-contacts}
 p_{ij}=x_{h_1}(\theta_{ij}),\quad
 q_{ij}=v_{ij}+x_{h_2}(\theta_{ij}+\pi),\quad
 g_{ij}=F_{ij}(\theta_{ij}),
\end{equation}
the equation $F_{ij}'=0$ gives $q_{ij}-p_{ij}=g_{ij}n(\theta_{ij})$.
The two supporting half-planes show that all cross-body distances are at
least $g_{ij}$, so these points realize the closest pair.  The clearance
bound excludes a third impact on this segment.  Positive curvature,
positive gap and the strict clearance margins give its alternating
contact collar.  Compactness of the parameter set and the finite number
of channels give a common collar and, by the relative theorem, a common
positive-offset collar.  No finite-horizon assertion is used.

One can also choose gates without using an unknown contact point as an
observation.  Let $\vartheta_{ij}^0=\arg(4-8i,5-10j)$.  The center
perturbation is less than $3/40$ and $d_{ij}>6.32$, so
$|\vartheta_{ij}-\vartheta_{ij}^0|<0.012$.
The equation $F_{ij}'=0$ and
$\|h_1'\|+\|h_2'\|\le11/200$ give
$|\theta_{ij}-\vartheta_{ij}|<0.009$.  Hence the first contact is
within $0.062$ of $n(\vartheta_{ij}^0)$ and the second within $0.131$
of $(4-8i,5-10j)^t-\tfrac34 n(\vartheta_{ij}^0)$.
Fixed planar disks of radius $1/4$ about these reference points therefore
contain all the appropriate contacts with a uniform interior margin.
On each obstacle these gates are disjoint.  Shrinking the contact collars
inside them preserves the strict third-obstacle clearance.  These are
prescribed selected gates; their common reference coordinates in this
proof are not passed to the intrinsic observer as an inter-channel
registration.

Finally, Lemma~\ref{lem:v23-signature-symmetry} and absence of proper
obstacle symmetries make every required incidence transition rigid.
The two deck sums in \eqref{eq:v44-cycles} are $\mathsf e_1$ and
$\mathsf e_2$, and both cycles start with $e_{00}$.  The stated tree
reaches the only other obstacle orbit.  This checks each part of the
anchoring and spanning hypotheses in a single realization.
\end{proof}

\subsection{Recovering the channel frames from the recovered curves}

For a complete convex image $C$, let
\[
 z_k(C)=\frac1{2\pi}\int_0^{2\pi}h_C(\theta)e^{-ik\theta}\,d\theta,
 \qquad k\ge2.
\]
These coefficients do not change under translation.  They are derived
from the boundary image recovered by analytic continuation, not added
to the endpoint-law observations.  Identify the oriented plane with
$\mathbb C$, so a unit complex number denotes a rotation.

\begin{proposition}[Two-harmonic registration in the realized family]
\label{prop:v44-harmonic-registration}
Let $C_{e,1}$ and $C_{e,2}$ be the two complete boundary images recovered
in the frame of channel $e$, whose origin is its first contact.  Set
$e_*=e_{00}$ and $c_{e,a}=s(C_{e,a})$.  The Euclidean congruence placing
the first-type copy of edge $e$ onto that of $e_*$ is
\begin{equation}\label{eq:v44-fourier-registration}
 U_e=\frac{z_3(C_{e,1})z_2(C_{e_*,1})}
           {z_3(C_{e_*,1})z_2(C_{e,1})},\qquad
 A_e(y)=U_e y+c_{e_*,1}-U_e c_{e,1}.
\end{equation}
In exact data $|U_e|=1$.  If $b_e=A_e(c_{e,2})$, then the lattice
columns and Gram matrix in this one common frame are
\begin{equation}\label{eq:v44-explicit-recovery}
 V=(b_{00}-b_{10}\ \ b_{00}-b_{01}),\qquad G=V^TV.
\end{equation}
The redundant fourth channel satisfies
\begin{equation}\label{eq:v44-redundant-cycle}
 b_{00}-b_{11}=(b_{00}-b_{10})+(b_{00}-b_{01}).
\end{equation}
On this compact family, registration and the displayed lattice recovery
are locally Lipschitz in the support functions of the recovered complete
curve images, in the uniform norm.  This statement does not assert a
Lipschitz analytic-continuation map from contact jets or endpoint laws.
\end{proposition}

\begin{proof}
Write the actual edge placement as $y\mapsto R_e y+p_e$, with rotation
angle $\theta_e$.  Its first-type image is
$C_{e,1}=R_e^T(K_{h_1}-p_e)$.  Therefore
\[
 z_k(C_{e,1})=e^{ik\theta_e}z_k(K_{h_1}),\qquad
 z_2(K_{h_1})=\epsilon/2,\quad
 z_3(K_{h_1})=(\epsilon/8)e^{i\pi/4}.
\]
Taking the ratio for mode $3$ and dividing by that for mode $2$ yields
$e^{i(\theta_e-\theta_{e_*})}$.  This proves the rotation formula without
selecting a branch of an argument.  Translation covariance of
\eqref{eq:v44-center} proves the formula for $A_e$.  Indeed
$c_{e,1}=R_e^T(-p_e)$, and hence
\[
 A_e(y)=R_{e_*}^TR_e y+R_{e_*}^T(p_e-p_{e_*}),\qquad
 b_{ij}=R_{e_*}^T(v_{ij}-p_{e_*}).
\]
Taking differences proves \eqref{eq:v44-explicit-recovery} and
\eqref{eq:v44-redundant-cycle}.  These differences are precisely the
translation holonomies of the two cycles, since their marked displacement
matrix is $I$.  The copy $C_{e_*,1}$ and the placed copy $C_{e_*,2}$,
together with $V\mathbb Z^2$, recover every obstacle translate.

For stability, the maps $h\mapsto z_k$ and $h\mapsto s$ have uniform-norm
operator bounds $1$ and $2$.  On the stated family
$|z_2|\ge1/400$ and $|z_3|\ge1/1600$.  The quotients in
\eqref{eq:v44-fourier-registration} are therefore locally Lipschitz
under small uniform perturbations; normalize a perturbed $U_e$ by its
modulus if it is not exactly a rotation.  This normalization is smooth
near the unit circle.  The centers lie in a common bounded set, so the
translation formula, the differences in $V$, and $V^TV$ inherit local
Lipschitz bounds.  This proof concerns the registration stage after the
complete curve images have been obtained.
\end{proof}

\subsection{The observation-to-table reconstruction in this family}

\begin{corollary}[Eight signed laws and four onsets]
\label{cor:v44-eight-laws}
There is $d_0>0$, common to the family, such that the four marked gaps and
the two same-type signed limiting endpoint laws at offset $d_0$ for each
of the four channels determine the whole marked table and its Euclidean
lattice up to one simultaneous proper Euclidean motion.  The parameters
$T,w,\epsilon,\delta$, the support coefficients, the contact curvatures,
and the placements of different channel frames are not supplied.
The two actions at a channel are not both even.  In particular, replacing
signed endpoint laws by their pointwise sign-erased versions is not the
observation used in this corollary.
\end{corollary}

\begin{proof}
Choose $d_0$ inside the uniform positive-offset collar in
Theorem~\ref{thm:v44-realized-family}.  At each edge and each contact type,
form the four-density ratio of Theorem~\ref{thm:v26-density-inverse} on a
common positive interior square.  A nonzero anchor recovers $S_{e,b}$
without knowing $B_{e,b}$ or the normalization.  In particular, if
$a_{e,b}=S_{e,b}''(0)$, then
\[
 \gamma_e=\operatorname{arsinh}(g_e\sqrt{a_{e,0}a_{e,1}}),\qquad
 \kappa_{e,b}=\frac{\cosh\gamma_e\sqrt{a_{e,b}/a_{e,1-b}}-1}{g_e}.
\]
For every $n\ge3$ the smooth finite-remainder theorem and the
homogeneous isolation in Theorem~\ref{thm:v22-signed-rigidity} give
$q_{e,n}=M_{e,n}^{-1}(s_{e,n}-R_{e,n})$.  The geometric bounds above do
not replace that nonlinear inverse; they ensure its hypotheses hold
simultaneously at all eight contacts.  The recovered analytic germs
determine the complete images $C_{e,a}$, by
Lemma~\ref{lem:v23-analytic-germ-globalization}.  Now use
Proposition~\ref{prop:v44-harmonic-registration} to place all edges in
the $e_{00}$ frame and recover $V$ and both obstacle images.  This is also
the concrete common-frame implementation of
Theorem~\ref{thm:v26-single-offset-global}.

If both actions in a channel were even, reflecting both transverse
graphs would leave the two actions unchanged.  The uniqueness of the
contact-jet inverse would make both graph jets even.  Analyticity would
then make each whole obstacle invariant under reflection in its contact
normal, contradicting Theorem~\ref{thm:v44-realized-family}.  Thus the
construction is not reduced to the reflection-symmetric local case.
\end{proof}

For a fully specified member, take
\begin{equation}\label{eq:v44-concrete-member}
 T=\begin{pmatrix}1&1/200\\0&1\end{pmatrix},\qquad
 w=(1/200,-1/200)^t,\qquad
 \epsilon=1/125,\quad\delta=3/400.
\end{equation}
Its four center vectors are $(4.030,4.995)$, $(-3.970,4.995)$,
$(3.980,-5.005)$ and $(-4.020,-5.005)$ in the construction frame.
Solving the strictly concave scalar problems
\eqref{eq:v44-gap-function} gives the actual contacts and onsets; these
are not assumed to lie on the center directions.  The exact recovered
Gram matrix, which is independent of the unknown starting-frame rotation,
is
\begin{equation}\label{eq:v44-concrete-gram}
 G=\begin{pmatrix}64&2/5\\2/5&40001/400\end{pmatrix}.
\end{equation}
The finite diagnostic supplied with the source solves all four contact
equations and computes \eqref{eq:v44-fourier-registration} from the
edge-frame support functions.  It does not fit the matrix $G$ directly
or give the reconstruction routine the generating matrix $T$.
The numerical comparison with \eqref{eq:v44-concrete-gram} checks this
finite geometric computation, not the all-order boundary-law theorem.

\subsection{Persistence beyond a finite-dimensional model}

\begin{proposition}[Analytic neighborhoods of the realization]
\label{prop:v44-analytic-neighborhood}
The family above has a neighborhood under real-analytic support-function
perturbations in which the four selected channels, strict clearance,
absence of obstacle symmetry, signature-rigid spanning structure and
rank-two anchoring pair persist.  One may independently add to $h_1,h_2$
any real-analytic $r_1,r_2$ with $\|r_a\|_{C^2}\le10^{-6}$.
Consequently the intrinsic determination applies on an
infinite-dimensional analytic neighborhood, not only to tables with the
finite-dimensional shape and placement description in
\eqref{eq:v44-supports}--\eqref{eq:v44-lattice-family}.
On each compact analytic subfamily satisfying the regularity conventions
of the global physical theorem, its charged long-bridge reconstruction
conclusion applies with the same four channel labels.
\end{proposition}

\begin{proof}
Here $\|r\|_{C^2}=\max_{0\le k\le2}\|r^{(k)}\|_\infty$.
The perturbation in curvature radius is at most $2\cdot10^{-6}$ and
that in the enclosing radius at most $10^{-6}$.  All separation,
curvature and clearance inequalities retain strict margins.  The scalar
contact maximum remains nondegenerate, so the channels and their gate
margins persist.  Fourier coefficients change by at most $10^{-6}$.
In particular the nonzero modes $2$ and $3$, of minimum moduli $1/400$
and $1/1600$, remain nonzero and exclude every proper rotation.

For reflections, write the arguments of these two coefficients as
$\beta_2,\beta_3$.  A necessary condition is
$2\beta_3-3\beta_2\in\pi\mathbb Z$.  In the unperturbed family
this quantity is $\pi/2$ or $\pi/3$.  Each argument changes by at most
$2\cdot10^{-6}/(1/1600)\le0.0032$, so the required congruence remains
impossible.  The vector \eqref{eq:v44-center}, which may now move,
removes the translation part of any candidate symmetry.  The two
perimeters remain different.  Thus signature uniqueness persists, and
the same combinatorial cycles and tree prove the intrinsic determination.
The two-harmonic formula continues to hold because the same unknown
coefficients cancel between copies; it does not require the absence of
higher Fourier modes.

For example, fixing any strip width $\sigma>0$ permits infinitely many
independent sufficiently small perturbations $a_n\cos n\theta$ and
$b_n\sin n\theta$, with exponentially summable coefficients.  Hence this
is genuinely an analytic neighborhood rather than a finite-dimensional
reparametrization.  On a compact analytic subfamily the positive margins
and the hypotheses of the observable calibration and global physical
reconstruction theorems give a common collar, the finite separating tests,
and the charged preparation caps.  Running their budget-indexed stages
with the present four labels proves the last assertion.  Compactness is
used at this last step; the whole bounded analytic neighborhood is not
silently asserted to be compact in an unspecified topology.
\end{proof}

The example thus implements the entire intrinsic route: nonlinear
long-bridge laws, amplitude-free action inversion, both contact jets,
analytic images, unique registration, and an uncalibrated rank-two lattice.
Its explicit registration bounds do not supply an effective rate for the
preceding analytic continuation.  In the richer planar-position protocol,
the direct-position benchmark remains available as well.  The two routes
continue to use different retained information, even on this same
nonsymmetric periodic family.
